\documentclass[a4paper,10pt,oneside,openany,final,article]{memoir}
\input{common}
\settocdepth{chapter}
\usepackage{minted}
\usepackage{fontspec}
% \setromanfont{Source Serif Pro}
% \setsansfont{Source Sans Pro}
% \setmonofont{Source Code Pro}

\begin{document}
\title{Free Value Or Else}
\author{
  Steve Downey \small<\href{mailto:sdowney@gmail.com}{sdowney@gmail.com}> \\
  Corentin Jabot \small<\href{mailto:corentin.jabot@gmail.com}{corentin.jabot@gmail.com}> \\
}
\date{} %unused. Type date explicitly below.
\maketitle

\begin{flushright}
  \begin{tabular}{ll}
    Document \#: & D4270R1 \\
    Date: & \today \\
    Project: & Programming Language C++ \\
    Audience: & LEWG
  \end{tabular}
\end{flushright}

\begin{abstract}
  \tcode{value_or} has two defects, and this paper carries the fix for
  each. The member half, from \cite{P3413R0}, adds
  \tcode{value_or_construct} and \tcode{value_or_else} to
  \tcode{std::optional} and \tcode{std::expected}, so that the alternative
  is built only when it is needed. The free half adds free function
  overloads of \tcode{value_or} --- and the related \tcode{reference_or},
  \tcode{or_invoke}, and \tcode{or_construct} --- over any nullable type,
  such as \tcode{std::optional}, \tcode{std::expected}, and pointers to
  objects, computing the type of the result correctly, via
  \tcode{common_type} and \tcode{common_reference}, as the member
  \tcode{std::optional::value_or} could not. The two were separate papers
  until LEWG asked for one. They compose, and they are severable.
\end{abstract}

\tableofcontents*

\chapter*{Changes Since Last Version}

\begin{itemize}
\item \textbf{R1}
  \begin{itemize}
  \item Unified with \cite{P3413R0} at LEWG's request: Corentin Jabot joins as
    co-author, and the members \tcode{value_or_construct} and
    \tcode{value_or_else} are folded into the proposal, the wording, and the
    implementation, covering \tcode{optional<T\&>}, which P3413R0 deferred and
    which is now C++26.
  \item An rvalue nullable moves its payload when its dereference expression
    permits moving. R0 dereferenced an lvalue always, observing and never
    consuming; the models now split by dereference behavior, so
    \tcode{optional} and \tcode{expected} move, while raw and smart pointers
    copy.
  \item \tcode{nullable} is proposed as a named, public concept. R0 specified
    it as exposition-only and asked LEWG for a poll.
  \item The concept is stated on \tcode{const remove_reference_t<T>\&} rather
    than \tcode{const T}, so the const-observability requirement holds for the
    deduced \tcode{T} of an lvalue argument, where the const was silently
    discarded before.
  \item \tcode{reference_or}'s second dangling guard is stated on the
    dereference type rather than on \tcode{T\&}, where it could never fire.
  \item \tcode{reference_or} mandates both that its result is a reference and
    that an rvalue argument is a \tcode{borrowed_nullable}. Temporary owning
    wrappers are therefore rejected even when dereferencing them yields
    \tcode{T\&}.
  \item \tcode{or_invoke} forwards its invocable, agreeing with the
    \tcode{invoke_result_t<I>} its return type is computed from;
    \tcode{or_construct} gains the engaged-path \emph{Mandates} and rejects
    reference result types.
  \item The asks to LEWG are cut from nine to zero. The dangling and borrowing
    checks are \emph{Mandates}, the concept stays minimal, the names stand,
    and the header is \tcode{<optional>}.
  \item Wording rebased to \cite{N5054}, which carries
    \tcode{optional<T\&>}.
  \end{itemize}
\item \textbf{R0} Initial revision.
\end{itemize}

\chapter{Comparison table}

\section{Raw pointer with a manual null check}

The fundamental nullable type is the raw pointer, and the fundamental idiom
for providing a default is the conditional expression guarded by a null
check. The free function gives that pattern a name, removes the repetition of
the checked entity, and computes the type of the result correctly.

\begin{tabular}{ lr }
  \begin{minipage}[t]{0.45\columnwidth}
    \begin{minted}[fontsize=\small]{c++}
Cat* cat = find_cat("Fido");
return cat ? *cat : Cat{"Default"};
    \end{minted}
  \end{minipage}
  &
    \begin{minipage}[t]{0.45\columnwidth}
      \begin{minted}[fontsize=\small]{c++}
Cat* cat = find_cat("Fido");
return std::value_or(cat, Cat{"Default"});
      \end{minted}
    \end{minipage}
\end{tabular}

\section{The fallback that is built anyway}

The argument to \tcode{value_or} is an ordinary function argument,
constructed on every call and destroyed unused on every engaged one. The
member \tcode{value_or_construct} takes the fallback's constructor
arguments instead, and builds nothing when the optional is engaged;
\tcode{value_or_else} defers an arbitrary computation the same way.

\begin{tabular}{ lr }
  \begin{minipage}[t]{0.45\columnwidth}
    \begin{minted}[fontsize=\small]{c++}
std::optional<Config> o = load();
// Config built even when o is engaged
return o.value_or(Config{"/etc/default"});
    \end{minted}
  \end{minipage}
  &
    \begin{minipage}[t]{0.45\columnwidth}
      \begin{minted}[fontsize=\small]{c++}
std::optional<Config> o = load();
// built only when o is disengaged
return o.value_or_construct("/etc/default");
      \end{minted}
    \end{minipage}
\end{tabular}

\section{The \tcode{value_or} truncation pitfall}

Because \tcode{std::optional<T>::value_or} converts its argument to
\tcode{T}, a perfectly reasonable-looking sentinel silently truncates. The
free function computes the common type of the contained value and the
alternative, and converts both \emph{outward} to that type instead of
forcing the alternative \emph{inward} through \tcode{T}.

\begin{tabular}{ lr }
  \begin{minipage}[t]{0.45\columnwidth}
    \begin{minted}[fontsize=\small]{c++}
std::optional<std::uint16_t> o = get();
auto v = o.value_or(-1);
// v is uint16_t{65535}
    \end{minted}
  \end{minipage}
  &
    \begin{minipage}[t]{0.45\columnwidth}
      \begin{minted}[fontsize=\small]{c++}
std::optional<std::uint16_t> o = get();
auto v = std::value_or(o, -1);
// v is int; -1 is preserved
      \end{minted}
    \end{minipage}
\end{tabular}

\section{Reference alternative without copies}

When both the contained value and the alternative are references to
long-lived objects, the result should be a reference. Member
\tcode{value_or} cannot do this; \tcode{reference_or} can, and statically
rejects the cases where the result would dangle.

\begin{tabular}{ lr }
  \begin{minipage}[t]{0.45\columnwidth}
    \begin{minted}[fontsize=\small]{c++}
std::optional<Logger&> opt = get_logger();
Logger& l =
    opt ? *opt : default_logger();
    \end{minted}
  \end{minipage}
  &
    \begin{minipage}[t]{0.45\columnwidth}
      \begin{minted}[fontsize=\small]{c++}
std::optional<Logger&> opt = get_logger();
Logger& l = std::reference_or(
    opt, default_logger());
      \end{minted}
    \end{minipage}
\end{tabular}

\chapter{Motivation}

\tcode{std::optional<T>::value_or} answers the most common question asked of
an optional --- ``the value, or this instead'' --- and gets two things
wrong. It answers with the wrong type, funneling both the contained value
and the alternative through \tcode{T}. And it pays for the alternative
whether or not it is wanted, because the alternative is an ordinary function
argument. The two defects are independent, they have been pursued
independently, and this paper carries the fix for each.

\section{The type of the answer}

\tcode{value_or} returns \tcode{T}, and it requires the alternative to be
convertible to \tcode{T}. Every use therefore funnels both the contained
value and the alternative through \tcode{T}, whether or not \tcode{T} is the
right type for the result. The consequences are familiar: sentinels truncate
(\tcode{optional<uint16_t>{}.value_or(-1)} is \tcode{65535}), wider
alternatives narrow, and reference results are simply impossible, so
\tcode{value_or} forces a copy even when both candidate results are
references to long-lived objects. It answers the wrong question correctly.

This is not fixable in place. Changing the return type of a member function
that has been shipping since C++17 is both an API and an ABI break, and the
\&\&-qualified overload set compounds the problem. During the work on
\tcode{std::optional<T\&>} \cite{P2988R12} the question was reopened in
earnest, because for an optional reference the cost of the wrong answer is
no longer a pessimization but a wrong program: returning \tcode{T} from
\tcode{optional<T\&>::value_or} copies, returning \tcode{T\&} dangles when
the alternative is a temporary. After extensive discussion in LEWG, P2988
concluded that ``there is no particularly wonderful solution for
\tcode{value_or} that does not involve a time machine,'' adopted the least
objectionable member behavior --- return \tcode{remove_cv_t<T>} by value ---
and explicitly forecast this paper: free functions over all types modeling
an optional-like concept, where the return type can finally be computed
correctly.

Existing practice diverges because there is no one member-function answer
that satisfies everyone:

\begin{center}
\begin{tabular}{ l l }
\hline
Implementation & Behavior \\
\hline
Standard & \tcode{optional<T>::value_or} returns a \tcode{T} \\
Boost & \tcode{optional<T>::value_or} returns a \tcode{T} \\
      & \tcode{optional<T\&>::value_or} returns a \tcode{T\&} \\
TartanLlama & \tcode{optional<T>::value_or} returns a \tcode{T} \\
      & \tcode{optional<T\&>::value_or} returns a \tcode{T} \\
Flux & returns the type of the ternary, similar to \tcode{common_reference} \\
Think-Cell & returns \tcode{common_reference}, with some caveats \\
\hline
\end{tabular}
\end{center}

The tools to compute the right answer postdate \tcode{optional}.
\tcode{std::common_type} existed but was not used;
\tcode{std::common_reference} arrived with ranges in C++20; and
\tcode{std::reference_constructs_from_temporary_v} \cite{P2255R2} arrived in
C++23, making it possible for the first time to \emph{statically reject}
the dangling cases that made reference-returning \tcode{value_or}
untenable. The conditional expression \tcode{b ? *m : u} has computed the
common type all along; the library can now name that computation, add the
lifetime check the core language does not perform, and do so once, for every
nullable type.

That last point is the generalization this paper takes from
\cite{P1255R13}: being ``contextually convertible to \tcode{bool} and
dereferenceable'' is a protocol, not a property of \tcode{std::optional}.
Raw pointers, \tcode{std::shared_ptr}, \tcode{std::unique_ptr},
\tcode{std::expected<T,E>}, and \tcode{std::optional<T\&>} all speak it.
Each of them today requires either a hand-written conditional or a
type-specific spelling of ``the value, or this instead.'' A single
constrained free function family covers them all, with one set of semantics
to learn and one place to get the corner cases right.

\section{The cost of the alternative}

The second defect is independent of the first, and older in the committee's
record. \tcode{o.value_or(Config{"/etc/default"})} constructs a
\tcode{Config} on every call and destroys it unused on every engaged one.
The parameter is an ordinary function parameter, so the alternative is
evaluated whether or not it is reached, and the cost is unbounded: an
expensive construction, an allocation, a lookup, a computation whose result
is thrown away. The workaround is the conditional expression the member was
supposed to replace, and callers who care about the cost write it.

The cure has been proposed before. \cite{P2218R0}, by Marc Mutz, proposed
\tcode{value_or_construct}, taking the fallback's constructor arguments
rather than a constructed fallback, so that nothing is built on the engaged
path; it did not proceed at the time. \cite{P3413R0} revived it, added
\tcode{value_or_else} for the general case of a lazily evaluated nullary
callable, and specified both on \tcode{std::optional} and
\tcode{std::expected}. Those members are the member half of this paper, and
their wording is reproduced below.

Laziness is orthogonal to the type of the result. \tcode{value_or_construct}
and \tcode{value_or_else} still return \tcode{T}, and still require the
alternative to converge on \tcode{T}. They fix what they set out to fix, on
the two types that have the member to grow.

\section{Two defects, two shapes}

The fixes take different shapes because the defects are of different kinds.
The cost is a defect in the parameters, and a new member name fixes it in
place: \tcode{value_or_construct} and \tcode{value_or_else} sit beside
\tcode{value_or} on \tcode{optional} and \tcode{expected}, spelled the way
the member they join is spelled, and nothing existing changes meaning. The
type is a defect in the signature, and no member can fix it. The return type
of a member shipping since C++17 can not change. A corrected second member
would have to be written, and kept in sync by hand, on every class that has
members at all. A raw pointer has none (see Why free functions, below).

So this paper proposes each defect's fix in the shape that defect admits: a
pair of members, and a family of free functions over the protocol. Neither
forecloses the other, and neither depends on the other.

\chapter{Design}

The proposal has two halves. The free half is four function templates over a
\tcode{nullable} concept:

\begin{minted}[fontsize=\small]{c++}
template <class T>
concept nullable = requires(const std::remove_reference_t<T>& t) {
    bool(t);
    *(t);
};

template <class T>
inline constexpr bool enable_borrowed_nullable = false;

template <class T>
concept borrowed_nullable =
  nullable<T> &&
  (std::is_lvalue_reference_v<T> ||
   enable_borrowed_nullable<std::remove_cvref_t<T>>);

template <class T>
inline constexpr bool enable_borrowed_nullable<T*> = true;

template <class T>
inline constexpr bool
  enable_borrowed_nullable<std::optional<T&>> = true;

template <class T, class E>
inline constexpr bool
  enable_borrowed_nullable<std::expected<T&, E>> = true;

// The type *m yields, carrying the value category of the nullable through
// to its payload.  Exposition only in the wording.
template <class T>
using deref_t = decltype(*std::declval<T>());

template <nullable T, class U,
          class R = std::common_reference_t<deref_t<T>, U&&>>
constexpr auto reference_or(T&& m, U&& u) -> R;

template <nullable T, class U,
          class R = std::common_type_t<deref_t<T>, U&&>>
constexpr auto value_or(T&& m, U&& u) -> R;

template <nullable T, class I,
          class R = std::common_type_t<deref_t<T>,
                                       std::invoke_result_t<I>>>
constexpr auto or_invoke(T&& m, I&& invocable) -> R;

template <class Ret = void, nullable T, class... Args,
          class R = std::conditional_t<std::is_void_v<Ret>,
                        std::remove_cvref_t<deref_t<T>>, Ret>>
constexpr R or_construct(T&& m, Args&&... args);

template <class Ret = void, nullable T, class E, class... Args,
          class R = std::conditional_t<std::is_void_v<Ret>,
                        std::remove_cvref_t<deref_t<T>>, Ret>>
constexpr R or_construct(T&& m, std::initializer_list<E> il, Args&&... args);
\end{minted}

All four have the same engaged-path behavior, producing
\tcode{*std::forward<T>(m)}
converted to \tcode{R}, and differ in how the alternative is supplied
(a value, a value bound as a reference, a nullary invocable, constructor
arguments) and in how \tcode{R} is computed (\tcode{common_type},
\tcode{common_reference}, or the element type itself).

The member half is \tcode{value_or_construct} and \tcode{value_or_else} on
\tcode{std::optional} and \tcode{std::expected}, as \cite{P3413R0}
specified them:

\begin{minted}[fontsize=\small]{c++}
template <class... Args>
constexpr T value_or_construct(Args&&... args) const&;

template <class U, class... Args>
constexpr T value_or_construct(std::initializer_list<U> il,
                               Args&&... args) const&;

template <class F>
constexpr T value_or_else(F&& f) const&;

// and the corresponding && overloads, on optional and expected alike
\end{minted}

Both return \tcode{T}, as the member \tcode{value_or} beside them does, and
neither constructs anything on the engaged path. The sections that follow
give the design of the free half and the choices made; the member
half is taken up under The member half, below.

\section{The \tcode{nullable} concept}

The concept is the ``maybe protocol'' of \cite{P1255R13}: an object is
nullable if it is contextually convertible to \tcode{bool} and
dereferenceable. The requirements are stated on a const lvalue: observation
must not require mutable access. The reference is stripped before the const
is applied, so the rule holds for the deduced \tcode{T} of an lvalue argument
as much as for an rvalue; spelled \tcode{const T} it would not, since const
applied to a reference type is discarded. A type whose \tcode{operator bool}
or \tcode{operator*} is non-const does not model \tcode{nullable}. Asking
``the value or a default'' is a query. A type that can not be queried
through const access is not answering it.

Types verified to model the concept: \tcode{std::optional<T>},
\tcode{std::expected<T,E>}, raw object pointers \tcode{T*},
\tcode{std::shared_ptr<T>}, \tcode{std::unique_ptr<T>}, and
\tcode{std::optional<T\&>} as specified by \cite{P2988R12}. Types verified
not to model it: arithmetic types and \tcode{bool} (contextually convertible
to \tcode{bool} but not dereferenceable), \tcode{std::string} and the
containers (neither), \tcode{std::nullptr_t}, \tcode{void}, and types
providing only one of the two operations. Iterators in general do not model
\tcode{nullable} because they are not contextually convertible to
\tcode{bool}. Correctly so: an iterator's validity is not a property the
iterator itself can report.

One possible refinement is deliberately not taken. P1255's
\tcode{nullable_object} also
requires that the result of dereferencing be an equality-preserving
\emph{object}, which excludes function pointers; the minimal concept above
admits them (a function pointer is contextually convertible to \tcode{bool}
and dereferenceable), and they are then rejected only downstream, when
\tcode{common_type} fails. We keep the concept minimal. Requiring
\tcode{is_object_v} of the referent would move that rejection to the
constraint, which reads better in a diagnostic, at the cost of a concept
users can no longer reproduce from its one-line statement.

\section{Borrowed nullables}

\tcode{nullable} says that an object can report whether dereferencing it is
valid. It does not say that the result of dereferencing survives destruction
of the nullable itself. That distinction does not matter to the three
value-producing functions, but it is decisive for \tcode{reference_or}.

A nullable can be understood as a view of zero or one element. In a
\tcode{borrowed_nullable}, it is that element --- the result of dereferencing
--- that is borrowed: its lifetime is managed independently of the nullable
wrapper. The name does not mean that the nullable object itself is borrowed.

This is the same lifetime distinction made by ranges. A
\tcode{single_view<T>} contains its element, and destroying the view
invalidates references to that element; an \tcode{optional<T>} has the same
relationship to its value. A \tcode{span<T>} presents externally owned
elements, as a raw pointer, \tcode{optional<T\&>}, or
\tcode{expected<T\&, E>} presents an externally owned referent.
\tcode{borrowed_range} permits an iterator or reference obtained from an
rvalue range to escape the range wrapper; \tcode{borrowed_nullable} permits
the reference obtained by dereferencing an rvalue nullable to escape through
\tcode{reference_or}.

As with \tcode{borrowed_range}, an lvalue nullable is always safe to observe:
the caller's object remains alive after the call. An rvalue must opt in via
\tcode{enable_borrowed_nullable}. Raw pointers opt in because destroying a
pointer does not destroy its pointee. \tcode{optional<T\&>} and
\tcode{expected<T\&, E>} opt in because they store references and destroying
the wrapper does not destroy the referred-to object.

Owning \tcode{optional} and \tcode{expected} do not opt in. Neither do
\tcode{unique_ptr} or \tcode{shared_ptr}: although their dereference operators
return \tcode{T\&}, a temporary \tcode{unique_ptr}, or a temporary
\tcode{shared_ptr} holding the last ownership share, destroys the referent at
the end of the full-expression. A custom deleter that happens not to delete is
not the lifetime guarantee on which a generic reference-returning function
may rely. Users may still pass every one of these types as an lvalue.

This rule has a deliberate ergonomic cost for an owning nullable returned by
value. Even if the selected reference would be consumed immediately, the
temporary wrapper cannot promise that its referent survives the return:

\begin{minted}[fontsize=\small]{c++}
reference_or(lookup(key), fallback); // ill-formed, even if consumed immediately

auto result = lookup(key);
reference_or(result, fallback);      // the wrapper remains alive

value_or(lookup(key), fallback);     // produce a value instead
\end{minted}

This is deliberately an opt-in promise. A program-defined nullable may
specialize \tcode{enable_borrowed_nullable} to \tcode{true} when destruction
of the nullable cannot invalidate the result of dereferencing it.

\section{Exposition-only, or a named public concept?}

We propose \tcode{std::nullable} as a named, public concept, usable outside
the standard library.

The protocol is vocabulary.
\cite{P2988R12} forecast these functions ``over all types modeling
optional-like, \tcode{concept std::maybe}'', and \cite{P1255R13} is titled,
in part, ``a concept to constrain maybes'' --- the concept has been treated
as a deliverable in its own right throughout this line of work. Users
writing their own functions over nullable types (a \tcode{value_or}
variant with different policy, an algorithm taking a range of nullables)
need the same constraint, and without a public name each codebase
re-derives it, and not always the same way. A named concept can be tested
directly (\tcode{static_assert(std::nullable<T>)}), appears by name in
diagnostics, and fixes the meaning of ``optional-like'' across the
ecosystem once.

There is a principled reason not to. A named standard concept is forever, and
a public name can never be demoted or meaningfully tightened, while an
exposition-only one can be promoted later without breaking anything. The
library has a directly analogous precedent: \tcode{\placeholder{boolean-testable}},
among the most-used constraints in the standard, is exposition-only for that
reason. However, the asymmetry cuts the other way here. A
constraint nobody can name is one every codebase re-derives, and the
re-derivations are what diverge; \tcode{\placeholder{boolean-testable}} is
plumbing for other concepts, where \tcode{nullable} is the thing users want
to write down. With \cite{P1255R13} concluded (see the discussion of
\tcode{yield_if} below), this paper is the concept's sole remaining home, so
there is no cross-paper coordination to wait on.

The semantics a public name has to carry are stated with it: the contextual
conversion is equality-preserving and does not modify the object; \tcode{*t}
may be evaluated only when that conversion yields \tcode{true}; and both
are valid on a const lvalue. Making the concept exposition-only instead would
cost nothing but the spelling, but this paper proposes the public name.

\section{Return type computation}

\tcode{deref_t<T>} is \tcode{decltype(*declval<T>())}, the type
\tcode{*m} yields, with the value category of the nullable carried through.
For a nullable passed as an lvalue with an \tcode{int} payload it is
\tcode{int\&}; for a \tcode{const optional<int>} it is \tcode{const
int\&}; for \tcode{const int*} it is \tcode{const int\&}; and for an
rvalue \tcode{optional<int>} it is \tcode{int\&\&}, on which see the next
section. The return type is then the common type
(for \tcode{value_or} and \tcode{or_invoke}) or the common reference (for
\tcode{reference_or}) of that dereference type and the alternative.

For \tcode{value_or}, \tcode{common_type} decays, so the result is a
prvalue, and mixed-type uses promote instead of truncating:

\begin{center}
\begin{tabular}{ l l l }
\hline
nullable \tcode{T} & alternative \tcode{U} & result \tcode{R} \\
\hline
\tcode{optional<int>} & \tcode{int} & \tcode{int} \\
\tcode{expected<int,E>}, \tcode{int*}, \tcode{shared_ptr<int>},
  \tcode{unique_ptr<int>} & \tcode{int} & \tcode{int} \\
\tcode{optional<int>} & \tcode{long} & \tcode{long} \\
\tcode{optional<int>} & \tcode{double} & \tcode{double} \\
\tcode{optional<uint16_t>} & \tcode{int} (e.g.\ \tcode{-1}) & \tcode{int} \\
\tcode{optional<string>} & \tcode{string} & \tcode{string} \\
\hline
\end{tabular}
\end{center}

This is the type the equivalent conditional expression would have produced,
which is the point: \tcode{value_or(m, u)} is the library spelling of
\tcode{m ? *m : u}, not of \tcode{m ? *m : T(u)}.

\section{Value categories, and what an rvalue nullable means}

The value category of the nullable reaches its payload. \tcode{deref_t<T>}
is \tcode{decltype(*declval<T>())}, so an expiring nullable dereferences to
an expiring payload, and the engaged path moves rather than copies:
\tcode{value_or(std::move(opt), u)} matches
\tcode{std::move(opt).value_or(u)}. The free functions do not lose the
member's efficiency.

That divides the behavior of the value-producing functions. The division is
not an accident of specification.
\tcode{std::optional} and
\tcode{std::expected} contain their value. An rvalue of one is an expiring
object whose payload expires with it, and \tcode{operator*} is
ref-qualified to say so. A pointer contains nothing. \tcode{T*},
\tcode{std::shared_ptr}, and \tcode{std::unique_ptr} dereference to
\tcode{T\&} even as rvalues, and therefore the value-producing functions copy
their referents. This rule follows the expression's value category; it makes
no independent ownership inference.

\tcode{std::optional<T\&>} settles which property is doing the work. It is
an \tcode{optional}, and it dereferences to \tcode{T\&}. Each nullable
reports what its own value category permits the value-producing functions to
do with its payload.

The alternative is to compute the dereference type from \tcode{T\&}
always, as \tcode{std::iter_reference_t} does, so that the functions
observe and never consume. It buys a rule that fits in a sentence: the
result is what the lvalue conditional expression would have produced. And
two smaller things. The return type of the value-producing functions stops
depending on the value category of the nullable, and there is no moved-from
path to reason about. However, it
buys them by making every owning nullable
copy a payload that is going away, which is the wrong answer stated
uniformly. We would rather be correct and pay the ergonomics.

The behavior is pinned either way. The engaged path is copy-counted for
\tcode{optional}, \tcode{expected}, and both smart pointers, so the
dereference result shows up as test output and a change of rule shows up as a
test failure (see Implementation Experience). This does not govern
\tcode{reference_or}; that operation separately requires
\tcode{borrowed_nullable}.

\section{\tcode{reference_or}}

\tcode{reference_or} accepts a \tcode{nullable} and mandates that it is a
\tcode{borrowed_nullable}, computes
\tcode{common_reference_t<deref_t<T>, U\&\&>}, and binds the result to the
contained value or to the alternative without copying either. Const
qualification flows through the
computation as it should: a \tcode{const optional<int>} or a \tcode{const
int\&} alternative each force \tcode{const int\&} as the result; an
\tcode{int*} with an \tcode{int\&} alternative yields \tcode{int\&},
writable through, and mutation through the result reaches the original
object. The function introduces no copies.

The common reference of two unrelated
lvalue references need not be a reference.
\tcode{common_reference_t<int\&, long\&>} is \tcode{long}, and
\tcode{common_reference_t<string\&, const char(\&)[2]>} is \tcode{string};
in such calls the common-reference computation creates a result object rather
than selecting either source object. \tcode{reference_or} therefore mandates
that its result type is a reference. Calls whose common reference is a value
are ill-formed and should use \tcode{value_or}; a function named for reference
semantics does not silently copy.

The danger of returning references is binding one to a temporary, and that
is now statically detectable. The implementation carries two guards:

\begin{minted}[fontsize=\small]{c++}
static_assert(!std::reference_constructs_from_temporary_v<R, U>);
static_assert(
    !std::reference_constructs_from_temporary_v<R, deref_t<T>>);
\end{minted}

These checks complement, rather than replace, \tcode{borrowed_nullable}.
The first rejects the alternative materializing a temporary: in
\tcode{reference_or(opt, 0)} the computed common reference of \tcode{int\&}
and \tcode{int\&\&} is \tcode{const int\&}, which would bind to a temporary
materialized from the literal, and the call is ill-formed. The second
rejects the symmetric case where the conversion of \tcode{*m} to \tcode{R}
would materialize a temporary. The mandate rejects an rvalue wrapper whose
destruction may invalidate an otherwise direct reference. Together the three
checks reject the known ways this operation could return a dangling
reference, including temporary owning smart pointers whose \tcode{operator*}
itself only reveals \tcode{T\&}.

These checks are specified as \emph{Mandates}, implemented by
\tcode{static_assert}s. A dangling call is therefore a hard,
clearly-diagnosed error wherever it occurs, including inside an
overload-resolution context, where it is an error and not a removed
candidate. As constraints, the overload disappears SFINAE-style, which
composes better with generic code probing for usability but turns a
lifetime bug into a silent ``no matching function'' or, worse, a fallback
to a different overload. Our position is that a dangling \tcode{reference_or}
call is a bug, and should be reported as one.

The Tokyo discussion of the P2988 homework concluded that \tcode{value_or}
is not an available name for a reference-returning function, and suggested
\tcode{ref_or}. This paper spells it \tcode{reference_or}, matching the
library's general preference for unabbreviated names.

\section{\tcode{or_invoke} and laziness}

\tcode{value_or} and \tcode{reference_or} evaluate their alternative
unconditionally. When it is expensive to construct, that cost is unwelcome
on the engaged path, and the established cure is to pass a nullary invocable
that is called only when needed. \tcode{or_invoke(m, f)} returns
\tcode{*m} converted to \tcode{R} when \tcode{m} is engaged and
\tcode{f()} converted to \tcode{R} otherwise, with
\tcode{R = common_type_t<deref_t<T>, invoke_result_t<I>>}; the
invocable is not called on the engaged path.

Because \tcode{or_invoke} uses \tcode{common_type}, an invocable returning
a reference still produces a value: \tcode{common_type_t<int\&, int\&>} is
\tcode{int}. The family therefore covers three of four quadrants, eager
value (\tcode{value_or}), eager reference (\tcode{reference_or}), and lazy
value (\tcode{or_invoke}, and \tcode{or_construct} below), and omits the
fourth. A lazy reference-returning form is specifiable (the dangling guard
would apply to \tcode{invoke_result_t<I>} just as it does to \tcode{U}), but
the use case, a reference alternative expensive to \emph{name} and not to
construct, is thin, and it is not proposed. That quadrant could compose
naturally as a separate \tcode{reference_invoke} proposal later.

The obvious name, \tcode{or_else}, is taken:
\tcode{std::optional::or_else} returns an \tcode{optional}, re-wrapping the
result, and giving the same name to a value-producing free function would
make \tcode{x.or_else(f)} and \tcode{or_else(x, f)} mean different things.
The member half faces the same collision, and answers it as \cite{P3413R0}
did, with \tcode{value_or_else}. For the free function this paper proposes
\tcode{or_invoke}. The two names differ because the two functions differ:
the member returns \tcode{T}, the free function the common type.

\section{\tcode{or_construct}: the fallback as constructor arguments}

\tcode{or_construct} is the free-function analogue of the other half's member
\tcode{value_or_construct}: the alternative is supplied not as a
value but as constructor \emph{arguments}, and the fallback object is
constructed lazily, only on the disengaged path. On the engaged path nothing
is constructed from them at all. An \tcode{initializer_list} overload
provides the same convenience the member form has.

However, its return-type rule differs from the other three. There is no
fallback \emph{value} whose type could participate in a
\tcode{common_type} computation, only arguments waiting to become one,
so by default \tcode{R} is the decayed element type itself,
\tcode{remove_cvref_t<deref_t<T>>}, and the conversion runs
\emph{inward}: the arguments construct an \tcode{R} directly, and the
engaged payload converts to \tcode{R}. This is the member function's model,
and here it is the right one; the outward common-type computation is
meaningless without a second value.

What the free function can offer that the member cannot is choice of the
constructed type. The leading \tcode{Ret} template parameter (defaulted
through a \tcode{void} sentinel) lets the caller name the result type
explicitly:

\begin{minted}[fontsize=\small]{c++}
std::optional<std::pmr::string> m = lookup(key);
auto s = std::or_construct<std::string>(m, "default");
// std::string either way: converted from the payload, or
// constructed from the arguments --- chosen by the caller
\end{minted}

The fallback is direct-initialized, so \tcode{explicit} constructors are
usable; the arguments were always explicitly a construction request. Left to
its default, \tcode{R} is a decayed object type, so \tcode{or_construct}
returns a prvalue and can not dangle. With \tcode{optional<T\&>} the result
is a value copy of the referent or a freshly constructed object, never a
reference. An explicitly named \tcode{Ret} is likewise required not to be a
reference, because otherwise fallback construction could bind the returned
reference to a temporary.

\section{The member half: \tcode{value_or_construct} and \tcode{value_or_else}}

The members are as \cite{P3413R0} specified them, and their design questions
are settled by the member they join. Both return \tcode{T}, and require the
alternative to converge on \tcode{T}, because a member named for
\tcode{value_or}, sitting beside it, would be a trap if it answered with a
different type; the corrected result type is the free half's business, under
its own names. Both come in \tcode{const\&} and \tcode{\&\&} form, so an
rvalue moves its payload out. \tcode{value_or_construct} takes constructor
arguments, with an \tcode{initializer_list} overload, and direct-initializes
the fallback, so \tcode{explicit} constructors are usable --- the arguments
were always explicitly a construction request. \tcode{value_or_else} takes a
nullary invocable whose result must convert to \tcode{T}. Neither touches
the fallback on the engaged path. \tcode{expected<void,E>} gains neither;
there is no value to construct.

Two things follow from the halves being one paper. The names do not collide:
\tcode{x.value_or_construct(args...)} and \tcode{or_construct<R>(m, args...)}
are different spellings of the same model, and \tcode{x.value_or_else(f)}
and \tcode{or_invoke(x, f)} differ in return type deliberately (see
\tcode{or_invoke} and laziness, above). And \tcode{optional<T\&>} no longer
has to be deferred: P3413 deferred it to \cite{P2988R12} while that paper
was in flight, and \tcode{optional<T\&>} is in C++26 now, so the members are
specified over it here, returning a value, exactly as its \tcode{value_or}
does. \tcode{expected<T\&,E>} is the one form still waiting on a paper of
its own, and the wording says so.

\section{Why free functions}

The member route is open for one defect and closed for the other, which is
why the proposal has two halves rather than one. Adding
\tcode{value_or_construct} and \tcode{value_or_else} adds new names, and new
names can be added; that is the member half. Correcting the type of the
answer can not be done that way. For \tcode{std::optional}
and \tcode{std::expected} the existing \tcode{value_or} cannot change its
return type, and a second member with corrected semantics would have to
coexist with the first forever, on each class, kept in sync by hand. For a
raw pointer there is nowhere to put a member at all, and adding one to each
smart pointer would repeat the same hand-synchronization. A constrained free
function template is the only form that can state the
semantics once and apply them to every type speaking the protocol,
including user-defined nullables, which model \tcode{nullable} without
doing anything to opt in.

These are ordinary function templates rather than customization point
objects. There is nothing to customize: the semantics are fully determined
by the protocol, and a type that wants different behavior for ``value or
default'' is not modeling \tcode{nullable}, it is doing something else. The usual
argument for a CPO, dispatching to a type's own implementation, is
therefore absent, and the names do not inhibit ADL.

\section{\tcode{std::expected}, and the error}

\tcode{std::expected<T,E>} models \tcode{nullable}: both \tcode{bool(e)}
and \tcode{*e} are valid on a const \tcode{expected}. Consequently
\tcode{value_or(e, u)} works, and discards the error, as the member
\tcode{value_or} does. This is by design: the moment the question is ``the
value, or this instead,'' the error has been decided to be uninteresting.
Uses that need the error have \tcode{error_or} and the monadic interface on
the member side; no error-observing overload is proposed here.

\section{What about \tcode{yield_if}, and the end of P1255}

\cite{P2988R12} forecast a fourth function, \tcode{yield_if}, producing a
view of zero or one elements, alongside the three proposed here. It was to
remain with the \cite{P1255R13} ranges work. However, that work is concluded.
Range support for \tcode{std::optional} \cite{P3168R2} supplanted
\tcode{views::maybe}, and \tcode{views::nullable} with it: an optional now
\emph{is} the view of zero or one elements, and a pointer-like nullable
reaches the same place through \tcode{optional<T\&>} \cite{P2988R12}, as
\tcode{p ? optional<T\&>(*p) : nullopt}. Downey, its author, does not expect
to revise P1255 further; this paper inherits its remaining deliverables, the
concept and the free functions over it.

That covers \tcode{yield_if} as well: \tcode{yield_if(cond, value)} is now
spellable as \tcode{cond ? optional(value) : nullopt}, which is itself a
range. What residual demand exists is for a named convenience in
comprehension-style pipelines; if implementation experience shows it is
wanted, it is a small separate paper. It is explicitly not part of this
proposal.

\chapter{P3413R0, and the two halves}

\cite{P3413R0} (``A more flexible \tcode{optional::value_or} (else!)''), by
Corentin Jabot, reviving \cite{P2218R0}, was a separate paper. It proposed
the member \tcode{value_or_construct}, constructing the fallback \tcode{T}
lazily and in place from forwarded arguments, with \tcode{initializer_list}
overloads, and \tcode{value_or_else}, taking a lazily evaluated nullary
callable, on both \tcode{std::optional} and \tcode{std::expected}. Reviewing
R0 of this paper in Brno, LEWG saw the two as fixes to the same member
function and asked, ``Can we ask you and Corentin to come back with one
proposal?'' This is that proposal. Corentin is a co-author here, and P3413's
members, wording, and rationale are the member half of this document; P3413R0
is not pursued separately.

Unification changed presentation, not design. The members are as P3413
specified them, with the one substantive addition the intervening time
allows: P3413 deferred \tcode{optional<T\&>} to \cite{P2988R12} while that
paper was in flight, and \tcode{optional<T\&>} is in C++26 now, so the
members are specified over it, and over the reference specializations of
\tcode{expected} as far as the working draft allows. The free half is
untouched by the unification; where it has changed since R0 --- the rvalue
rule of Design --- it changed for its own reasons.

What each half fixes stays distinct, and the distinction matters in review.
The member half fixes the \emph{cost} of the alternative: with
\tcode{value_or_construct} the fallback object is not constructed at all on
the engaged path, and \tcode{value_or_else} defers arbitrary computation. It
does not touch the return type; everything still returns \tcode{T}, with the
alternative still required to convert to \tcode{T}, so the truncation
pitfall, the forced copies, and the impossibility of reference results all
remain --- deliberately, since a member beside \tcode{value_or} has to
behave like \tcode{value_or}. It is also necessarily member-by-member:
\tcode{optional} and \tcode{expected} each grow six functions, and pointers
of any kind get nothing.

The free half fixes the \emph{type} of the answer and its \emph{reach}:
results are computed via
\tcode{common_type}/\tcode{common_reference} instead of funneled through
\tcode{T}, reference results become possible
and dangling-checked, and one set of free functions covers every nullable
type. The lazy cases are where the halves meet. \tcode{or_invoke} and member
\tcode{value_or_else} differ in form (free vs.\ member), reach (any
nullable vs.\ \tcode{optional}/\tcode{expected}), and return type
(\tcode{common_type} vs.\ \tcode{T}). \tcode{or_construct} and member
\tcode{value_or_construct} share the in-place, lazy construction model and
the \tcode{initializer_list} overload; the free function also lets
the caller name the constructed type
(\tcode{or_construct<string>(m, args...)}) and works on every nullable.

\begin{center}
\begin{tabular}{ l l l }
\hline
 & member half (P3413R0) & free half \\
\hline
form & members & free functions \\
applies to & \tcode{optional}, \tcode{expected} &
  any \tcode{nullable} (incl.\ pointers) \\
return type & \tcode{T} & \tcode{common_type} / \tcode{common_reference} \\
truncation pitfall & unchanged, by design & fixed \\
reference results & no & \tcode{reference_or}, dangling-checked \\
lazy alternative & \tcode{value_or_else} & \tcode{or_invoke} \\
in-place fallback construction & \tcode{value_or_construct} &
  \tcode{or_construct}, caller-chosen result type \\
rvalue nullable moves payload & yes (\tcode{\&\&} overloads) &
  yes, for the owning models (see Design) \\
\tcode{optional<T\&>} & members return a value & covered by the concept \\
\hline
\end{tabular}
\end{center}

The halves compose. Adopt the members and \tcode{optional} and
\tcode{expected} gain the lazy, in-place fallback, spelled the way a member
is spelled and moving out of an rvalue. Adopt the free functions and the
corrected semantics --- the computed result type, reference results, the
reach to pointers and every other nullable --- become available generically,
over those same types and over the ones that have no members to grow. A
caller writing \tcode{x.value_or_construct(args...)} and one writing
\tcode{or_construct<R>(m, args...)} are served by the same design at two
spellings. The halves also remain severable: LEWG can take one and leave the
other, and each stands alone. What unification bought is that they need no
longer be weighed against each other.

\chapter{Principles for Reification of Design}

The same principles that governed \cite{P2988R12} govern here.

\textbf{Errors that can be detected at compile time must be detected at
compile time.} The concept rejects non-nullable arguments at the
constraint; \tcode{reference_or} additionally mandates against rvalue
wrappers that do not model \tcode{borrowed_nullable}, non-reference
common-reference results, and every conversion that would bind the result to
a temporary.

\textbf{Never knowingly produce a dangling reference.} The functions take
no address and extend no lifetime; \tcode{reference_or}'s result always
refers to an object that existed before the call. The second guard is stated
on \tcode{deref_t<T>}, the type \tcode{*m} yields; stated on
\tcode{T\&} it would name the nullable rather than its payload, and could
never fire. The rejected cases are those where the core language would
have materialized a temporary: a prvalue alternative whose common reference
with the dereference type is a const reference
(\tcode{reference_or(opt, 0)}), and conversions of the contained value
to the result type that materialize.
The \tcode{borrowed_nullable} mandate covers the orthogonal danger that
destroying a temporary nullable destroys its directly referenced payload.

\textbf{Value-producing functions never dangle by construction.}
\tcode{value_or}, \tcode{or_invoke}, and \tcode{or_construct} return
non-reference values; \tcode{or_construct} mandates this even when its result
type is explicitly selected.

These principles are enforced by a negative-compilation test suite in the
reference implementation: each forbidden combination is a translation unit
that must fail to compile, with the diagnostic matched against the
constraint or \tcode{static_assert} that is required to fire.

\chapter{Implementation Experience}

The reference implementation is \tcode{beman.free_value_or}
\cite{Downey_free_value_or}. The implementation floor is C++23, required
for \tcode{reference_constructs_from_temporary_v}; the functions are
otherwise C++20 material.

The test suite exercises a matrix:

\begin{itemize}
\item the nullable types \tcode{std::optional<T>},
  \tcode{std::expected<T,E>}, \tcode{T*}, \tcode{std::shared_ptr<T>},
  \tcode{std::unique_ptr<T>}, and --- gated on a C++26 toolchain, via the
  Beman reference implementations --- \tcode{std::optional<T\&>} and
  \tcode{std::expected<T\&,E>};
\item engaged and disengaged;
\item the nullable as lvalue, const lvalue, and rvalue, the rvalue cases
  copy- and move-counted on the engaged path, so that dereference behavior is
  pinned rather than assumed;
\item the alternative as lvalue and rvalue, and mixed-type combinations;
\item constant evaluation of all four functions;
\item laziness of \tcode{or_invoke} and \tcode{or_construct}, and
  eagerness of \tcode{value_or};
\item for \tcode{or_construct}, construction arities from zero arguments
  through multi-argument emplacement and the \tcode{initializer_list}
  overload, with construction confirmed to occur only on the disengaged
  path; and
\item negative-compilation cases for concept violations, non-reference
  common-reference results, conversion-created temporaries, and temporary
  owning \tcode{optional}, \tcode{expected}, \tcode{unique_ptr}, and
  \tcode{shared_ptr} wrappers; and for \tcode{or_construct}, an explicitly
  selected reference result, a payload that cannot be explicitly converted,
  and fallback arguments that cannot construct the result.
\end{itemize}

Reference identity and mutation
round-trips confirm that \tcode{reference_or} introduces no copies:
writes through its result reach the contained object when engaged and the
alternative when not. The \tcode{optional<T\&>} suite additionally verifies
rebinding semantics: after an \tcode{optional<T\&>} is rebound to a new
referent, \tcode{value_or} produces and \tcode{reference_or} refers to the
new referent, and disengaging falls through to the alternative.

The suite found two behaviors reported in Design: that a non-reference
common-reference result must be rejected, and that \tcode{or_invoke}'s use of
\tcode{common_type} decays reference-returning invocables to values,
leaving the lazy-reference quadrant of the family empty by choice. The
negative-compilation harness was itself validated: each WILL-FAIL
translation unit was checked to fail for the required reason by mutating
its expected-diagnostic pattern and observing the test fail. The rvalue
cases were validated the same way, against the library and not the
harness. Reverting the dereference type to \tcode{std::iter_reference_t<T>},
the observe-only alternative, makes \tcode{reference_or}'s static assertions
fail to compile and, with those removed so that they cannot mask the rest,
fails the move counts for \tcode{value_or} over \tcode{optional} and over
\tcode{expected}, for \tcode{or_invoke}, and for \tcode{or_construct}. It
leaves the smart pointer cases passing, because those copy under both rules.
That is the dereference rule of Design showing up as test output. Coverage of the
header from that one file is complete on lines, functions, and branches.

Verified \tcode{common_reference} results for \tcode{reference_or},
showing const propagation:

\begin{center}
\begin{tabular}{ l l l }
\hline
nullable (lvalue) & alternative (lvalue) & result \\
\hline
\tcode{optional<int>\&} & \tcode{int\&} & \tcode{int\&} \\
\tcode{const optional<int>\&} & \tcode{int\&} & \tcode{const int\&} \\
\tcode{optional<int>\&} & \tcode{const int\&} & \tcode{const int\&} \\
\tcode{expected<int,int>\&} & \tcode{int\&} & \tcode{int\&} \\
\tcode{int*} & \tcode{int\&} & \tcode{int\&} \\
\tcode{shared_ptr<int>\&} & \tcode{int\&} & \tcode{int\&} \\
\tcode{optional<string>\&} & \tcode{string\&} & \tcode{string\&} \\
\hline
\end{tabular}
\end{center}

The current implementation, less the portability shims --- a polyfill over
the compiler builtin for \tcode{reference_constructs_from_temporary_v}, and a
template parameter order that MSVC accepts:

\begin{minted}[fontsize=\small]{c++}
template <class T>
concept nullable = requires(const std::remove_reference_t<T>& t) {
    bool(t);
    *(t);
};

template <class T>
inline constexpr bool enable_borrowed_nullable = false;

template <class T>
concept borrowed_nullable =
    nullable<T> && (std::is_lvalue_reference_v<T> ||
                    enable_borrowed_nullable<std::remove_cvref_t<T>>);

template <class T>
inline constexpr bool enable_borrowed_nullable<T*> = true;

template <class T>
inline constexpr bool enable_borrowed_nullable<std::optional<T&>> = true;

template <class T, class E>
inline constexpr bool enable_borrowed_nullable<std::expected<T&, E>> = true;

// The type *m yields, carrying the value category of the nullable through
// to its payload.  Exposition only in the wording.
template <class T>
using deref_t = decltype(*std::declval<T>());

template <nullable T, class U,
          class R = std::common_reference_t<deref_t<T>, U&&>>
constexpr auto reference_or(T&& m, U&& u) -> R {
    static_assert(borrowed_nullable<T>);
    static_assert(std::is_reference_v<R>);
    static_assert(!std::reference_constructs_from_temporary_v<R, U>);
    static_assert(
        !std::reference_constructs_from_temporary_v<R, deref_t<T>>);
    return bool(m) ? static_cast<R>(*std::forward<T>(m))
                   : static_cast<R>((U&&)u);
}

template <nullable T, class U,
          class R = std::common_type_t<deref_t<T>, U&&>>
constexpr auto value_or(T&& m, U&& u) -> R {
    return bool(m) ? static_cast<R>(*std::forward<T>(m))
                   : static_cast<R>(std::forward<U>(u));
}

template <nullable T, class I,
          class R = std::common_type_t<deref_t<T>,
                                       std::invoke_result_t<I>>>
constexpr auto or_invoke(T&& m, I&& invocable) -> R {
    return bool(m) ? static_cast<R>(*std::forward<T>(m))
                   : static_cast<R>(std::forward<I>(invocable)());
}

template <class Ret = void, nullable T, class... Args,
          class R = std::conditional_t<std::is_void_v<Ret>,
                        std::remove_cvref_t<deref_t<T>>, Ret>>
constexpr R or_construct(T&& m, Args&&... args) {
    static_assert(!std::is_reference_v<R>);
    static_assert(std::is_constructible_v<R, Args...>);
    static_assert(requires { static_cast<R>(*std::declval<T>()); });
    return bool(m) ? static_cast<R>(*std::forward<T>(m))
                   : R(std::forward<Args>(args)...);
}

template <class Ret = void, nullable T, class E, class... Args,
          class R = std::conditional_t<std::is_void_v<Ret>,
                        std::remove_cvref_t<deref_t<T>>, Ret>>
constexpr R or_construct(T&& m, std::initializer_list<E> il,
                         Args&&... args) {
    static_assert(!std::is_reference_v<R>);
    static_assert(
        std::is_constructible_v<R, std::initializer_list<E>&, Args...>);
    static_assert(requires { static_cast<R>(*std::declval<T>()); });
    return bool(m) ? static_cast<R>(*std::forward<T>(m))
                   : R(il, std::forward<Args>(args)...);
}
\end{minted}

The member side is implemented too. \tcode{value_or_construct} and
\tcode{value_or_else} were added to the Beman reference implementations of
\tcode{optional} \cite{The_Beman_Project_beman_optional} and \tcode{expected}
\cite{The_Beman_Project_beman_expected}, vendored into
\tcode{beman.free_value_or} and edited in place, on both the primary
templates and their reference specializations. \tcode{beman.expected} already
carries \tcode{expected<T\&,E>}, \tcode{expected<T,E\&>}, and
\tcode{expected<T\&,E\&>}, so the members run over references now, ahead of any
such paper reaching the working draft; standardizing those forms is the only
dependency, and not implementing them. The two libraries' own suites,
GoogleTest for \tcode{optional} and Catch2 for \tcode{expected}, run as
part of this project's tests, and cover laziness (a construction counter that never
ticks on the engaged path), the \tcode{initializer_list} overload, rvalue
move-out, constant evaluation, and the reference specializations. All green.

\chapter{Proposal}

Add to the standard library four constrained free function templates,
\tcode{value_or}, \tcode{reference_or}, \tcode{or_invoke}, and
\tcode{or_construct}, the named concept \tcode{nullable}, and the
\tcode{borrowed_nullable} refinement and customization point, with the
semantics described above, targeting C++29. The value-producing functions
should be usable on every type modeling \tcode{nullable}, including
\tcode{std::optional}, \tcode{std::expected}, raw pointers,
\tcode{std::shared_ptr}, \tcode{std::unique_ptr}, and
\tcode{std::optional<T\&>}. \tcode{reference_or} accepts all of them as
lvalues, and accepts rvalues only when they model \tcode{borrowed_nullable}.

Add the member functions of the other half, as \cite{P3413R0} specified
them:
\tcode{value_or_construct}, with its \tcode{initializer_list} overload,
and \tcode{value_or_else}, on both \tcode{std::optional} and
\tcode{std::expected}, each in \tcode{const\&} and \tcode{\&\&} form,
producing the fallback lazily and in place. Their wording is reproduced
below, unchanged in substance from P3413. \tcode{std::optional<T\&>}
\cite{P2988R12} is in C++26, and the members apply to it directly, returning a
value; a future \tcode{std::expected<T\&,E>} would carry them the same way, and
that one case depends on the expected-over-references work, as the wording
notes. The two feature sets are severable: LEWG may take the free functions,
the members, or both, but the point of unifying is that it need not.

\tcode{reference_or} mandates that its computed common reference is a
reference type. An rvalue nullable moves its payload in the value-producing
functions exactly when dereferencing that rvalue produces an rvalue
reference. Thus \tcode{optional} and \tcode{expected} move their values,
while raw and smart pointers copy their referents. Dereferencing an lvalue
always would make every owning nullable copy a payload that is going away;
this paper instead follows the rvalue's dereference expression (see Design).
\tcode{or_construct} rejects an explicitly selected reference result type.
The conversion dangling guards are \emph{Mandates}, so a
dangling call is diagnosed and not
removed from the overload set. \tcode{nullable} is a public concept.
The names are \tcode{value_or}, \tcode{reference_or}, \tcode{or_invoke},
and \tcode{or_construct}. The concept stays minimal, so a function pointer
models it and is rejected downstream when \tcode{common_type} fails, and not
at the constraint. The functions do not inhibit ADL. They should be
freestanding, and we propose \tcode{<optional>} as the header, where users
will look for them.

\chapter{Wording}

Initial draft wording, relative to \cite{N5054}, to be completed after LEWG
design review.

For a type \tcode{T}, let \exposid{deref-t}\tcode{<T>} denote
\tcode{decltype(*declval<T>())}.

Add to the \tcode{<optional>} header synopsis:

\begin{codeblock}
namespace std {
  template<class T>
    concept nullable = requires(const remove_reference_t<T>& t) {
      bool(t);
      *(t);
    };

  template<class T>
    inline constexpr bool enable_borrowed_nullable = false;

  template<class T>
    concept borrowed_nullable =
      nullable<T> && (is_lvalue_reference_v<T> ||
                      enable_borrowed_nullable<remove_cvref_t<T>>);

  template<class T>
    inline constexpr bool enable_borrowed_nullable<T*> = true;
  template<class T>
    inline constexpr bool enable_borrowed_nullable<optional<T&>> = true;

  template<nullable T, class U>
    constexpr common_reference_t<@\exposid{deref-t}@<T>, U&&>
      reference_or(T&& m, U&& u);

  template<nullable T, class U>
    constexpr common_type_t<@\exposid{deref-t}@<T>, U&&>
      value_or(T&& m, U&& u);

  template<nullable T, class I>
    constexpr common_type_t<@\exposid{deref-t}@<T>, invoke_result_t<I>>
      or_invoke(T&& m, I&& invocable);

  template<class Ret = void, nullable T, class... Args>
    constexpr conditional_t<is_void_v<Ret>,
                            remove_cvref_t<@\exposid{deref-t}@<T>>, Ret>
      or_construct(T&& m, Args&&... args);

  template<class Ret = void, nullable T, class E, class... Args>
    constexpr conditional_t<is_void_v<Ret>,
                            remove_cvref_t<@\exposid{deref-t}@<T>>, Ret>
      or_construct(T&& m, initializer_list<E> il, Args&&... args);
}
\end{codeblock}

Add to the \tcode{<expected>} header synopsis:

\begin{codeblock}
namespace std {
  template<class T, class E>
    inline constexpr bool enable_borrowed_nullable<expected<T&, E>> = true;
}
\end{codeblock}

The \tcode{expected<T\&, E>} specialization depends on standardizing
\tcode{expected} over reference types.

\pnum
Let \tcode{t} be a const lvalue of type \tcode{remove_reference_t<T>}.
\tcode{T} models \tcode{nullable} only if:

\begin{itemize}
\item \tcode{bool(t)} is equality-preserving and does not modify \tcode{t};
\item \tcode{*t} is valid only when \tcode{bool(t)} is \tcode{true}; and
\item \tcode{*t} is equality-preserving and does not modify \tcode{t}.
\end{itemize}

\pnum
The value of \tcode{enable_borrowed_nullable<T>} may be set to \tcode{true}
for a cv-unqualified program-defined type \tcode{T} for which destroying an
object of type \tcode{T} cannot invalidate the result of dereferencing that
object. Such a specialization shall be usable in constant expressions and
have type \tcode{const bool}.

\begin{itemdescr}
\begin{codeblock}
template<nullable T, class U>
  constexpr common_reference_t<@\exposid{deref-t}@<T>, U&&>
    reference_or(T&& m, U&& u);
\end{codeblock}

\pnum
Let \tcode{R} be \tcode{common_reference_t<@\exposid{deref-t}@<T>, U\&\&>}.

\pnum
\mandates
\tcode{borrowed_nullable<T>} is \tcode{true},
\tcode{is_reference_v<R>} is \tcode{true},
\tcode{reference_constructs_from_temporary_v<R, U>} is \tcode{false}, and
\tcode{reference_constructs_from_temporary_v<R, @\exposid{deref-t}@<T>>} is
\tcode{false}.

\pnum
\effects
Equivalent to:
\tcode{return bool(m) ? static_cast<R>(*std::forward<T>(m)) :}\newline
\tcode{\ \ \ \ \ \ \ static_cast<R>(std::forward<U>(u));}
\end{itemdescr}

\begin{itemdescr}
\begin{codeblock}
template<nullable T, class U>
  constexpr common_type_t<@\exposid{deref-t}@<T>, U&&>
    value_or(T&& m, U&& u);
\end{codeblock}

\pnum
Let \tcode{R} be \tcode{common_type_t<@\exposid{deref-t}@<T>, U\&\&>}.

\pnum
\effects
Equivalent to:
\tcode{return bool(m) ? static_cast<R>(*std::forward<T>(m)) :}\newline
\tcode{\ \ \ \ \ \ \ static_cast<R>(std::forward<U>(u));}
\end{itemdescr}

\begin{itemdescr}
\begin{codeblock}
template<nullable T, class I>
  constexpr common_type_t<@\exposid{deref-t}@<T>, invoke_result_t<I>>
    or_invoke(T&& m, I&& invocable);
\end{codeblock}

\pnum
Let \tcode{R} be
\tcode{common_type_t<@\exposid{deref-t}@<T>, invoke_result_t<I>>}.

\pnum
\effects
Equivalent to:
\tcode{return bool(m) ? static_cast<R>(*std::forward<T>(m)) :}\newline
\tcode{\ \ \ \ \ \ \ static_cast<R>(std::forward<I>(invocable)());}
\end{itemdescr}

\begin{itemdescr}
\begin{codeblock}
template<class Ret = void, nullable T, class... Args>
  constexpr conditional_t<is_void_v<Ret>,
                          remove_cvref_t<@\exposid{deref-t}@<T>>, Ret>
    or_construct(T&& m, Args&&... args);
\end{codeblock}

\pnum
Let \tcode{R} be \tcode{conditional_t<is_void_v<Ret>,
remove_cvref_t<@\exposid{deref-t}@<T>>, Ret>}.

\pnum
\mandates
\tcode{is_reference_v<R>} is \tcode{false},
\tcode{is_constructible_v<R, Args...>} is \tcode{true}, and the expression
\tcode{static_cast<R>(*declval<T>())} is well-formed.

\pnum
\effects
Equivalent to:
\tcode{return bool(m) ? static_cast<R>(*std::forward<T>(m)) :}\newline
\tcode{\ \ \ \ \ \ \ R(std::forward<Args>(args)...);}
\end{itemdescr}

\begin{itemdescr}
\begin{codeblock}
template<class Ret = void, nullable T, class E, class... Args>
  constexpr conditional_t<is_void_v<Ret>,
                          remove_cvref_t<@\exposid{deref-t}@<T>>, Ret>
    or_construct(T&& m, initializer_list<E> il, Args&&... args);
\end{codeblock}

\pnum
Let \tcode{R} be \tcode{conditional_t<is_void_v<Ret>,
remove_cvref_t<@\exposid{deref-t}@<T>>, Ret>}.

\pnum
\mandates
\tcode{is_reference_v<R>} is \tcode{false},
\tcode{is_constructible_v<R, initializer_list<E>\&, Args...>} is \tcode{true},
and the expression \tcode{static_cast<R>(*declval<T>())} is well-formed.

\pnum
\effects
Equivalent to:
\tcode{return bool(m) ? static_cast<R>(*std::forward<T>(m)) :}\newline
\tcode{\ \ \ \ \ \ \ R(il, std::forward<Args>(args)...);}
\end{itemdescr}

\section{Member functions, from P3413R0}

The following member wording is reproduced from \cite{P3413R0}, relative to
\cite{N5054}, and is proposed jointly with the free functions above. It is
purely additive: the existing \tcode{value_or} is unchanged. Return types and
\emph{Effects} follow P3413; the \emph{Mandates} spell the construction
requirement as \tcode{is_constructible_v<T, Args...>}, matching the
\emph{Effects}.

\subsection{\tcode{optional} [optional.optional], [optional.observe]}

To the synopsis of \tcode{optional} in [optional.optional.general], after the
\tcode{value_or} declarations, add:

\begin{codeblock}
template<class... Args> constexpr T value_or_construct(Args&&... args) const&;
template<class... Args> constexpr T value_or_construct(Args&&... args) &&;

template<class U, class... Args>
  constexpr T value_or_construct(initializer_list<U> il, Args&&... args) const&;
template<class U, class... Args>
  constexpr T value_or_construct(initializer_list<U> il, Args&&... args) &&;

template<class F> constexpr T value_or_else(F&& f) const&;
template<class F> constexpr T value_or_else(F&& f) &&;
\end{codeblock}

To [optional.observe], add:

\begin{itemdescr}
\begin{codeblock}
template<class... Args> constexpr T value_or_construct(Args&&... args) const&;
\end{codeblock}
\pnum
\mandates
\tcode{is_copy_constructible_v<T>} is \tcode{true} and
\tcode{is_constructible_v<T, Args...>} is \tcode{true}.
\pnum
\effects
Equivalent to:
\tcode{return has_value() ? **this : T(std::forward<Args>(args)...);}
\end{itemdescr}

\begin{itemdescr}
\begin{codeblock}
template<class... Args> constexpr T value_or_construct(Args&&... args) &&;
\end{codeblock}
\pnum
\mandates
\tcode{is_move_constructible_v<T>} is \tcode{true} and
\tcode{is_constructible_v<T, Args...>} is \tcode{true}.
\pnum
\effects
Equivalent to:
\tcode{return has_value() ? std::move(**this) : T(std::forward<Args>(args)...);}
\end{itemdescr}

\begin{itemdescr}
\begin{codeblock}
template<class U, class... Args>
  constexpr T value_or_construct(initializer_list<U> il, Args&&... args) const&;
\end{codeblock}
\pnum
\mandates
\tcode{is_copy_constructible_v<T>} is \tcode{true} and
\tcode{is_constructible_v<T, initializer_list<U>\&, Args...>} is \tcode{true}.
\pnum
\effects
Equivalent to:
\tcode{return has_value() ? **this : T(il, std::forward<Args>(args)...);}
\end{itemdescr}

\begin{itemdescr}
\begin{codeblock}
template<class U, class... Args>
  constexpr T value_or_construct(initializer_list<U> il, Args&&... args) &&;
\end{codeblock}
\pnum
\mandates
\tcode{is_move_constructible_v<T>} is \tcode{true} and
\tcode{is_constructible_v<T, initializer_list<U>\&, Args...>} is \tcode{true}.
\pnum
\effects
Equivalent to:
\tcode{return has_value() ? std::move(**this) : T(il, std::forward<Args>(args)...);}
\end{itemdescr}

\begin{itemdescr}
\begin{codeblock}
template<class F> constexpr T value_or_else(F&& f) const&;
\end{codeblock}
\pnum
Let \tcode{U} be \tcode{invoke_result_t<F>}.
\pnum
\mandates
\tcode{is_copy_constructible_v<T>} is \tcode{true} and
\tcode{is_convertible_v<U, T>} is \tcode{true}.
\pnum
\effects
Equivalent to:
\tcode{return has_value() ? **this : std::forward<F>(f)();}
\end{itemdescr}

\begin{itemdescr}
\begin{codeblock}
template<class F> constexpr T value_or_else(F&& f) &&;
\end{codeblock}
\pnum
Let \tcode{U} be \tcode{invoke_result_t<F>}.
\pnum
\mandates
\tcode{is_move_constructible_v<T>} is \tcode{true} and
\tcode{is_convertible_v<U, T>} is \tcode{true}.
\pnum
\effects
Equivalent to:
\tcode{return has_value() ? std::move(**this) : std::forward<F>(f)();}
\end{itemdescr}

\pnum
The reference specialization \tcode{optional<T\&>} \cite{P2988R12} carries
these members too; its wording is given together with
\tcode{expected<T\&,E>} below.

\subsection{\tcode{expected} [expected.expected], [expected.object.obs]}

To the synopsis of \tcode{expected} in [expected.expected.general], after the
\tcode{value_or} declarations, add:

\begin{codeblock}
template<class... Args> constexpr T value_or_construct(Args&&... args) const&;
template<class... Args> constexpr T value_or_construct(Args&&... args) &&;

template<class U, class... Args>
  constexpr T value_or_construct(initializer_list<U> il, Args&&... args) const&;
template<class U, class... Args>
  constexpr T value_or_construct(initializer_list<U> il, Args&&... args) &&;

template<class F> constexpr T value_or_else(F&& f) const&;
template<class F> constexpr T value_or_else(F&& f) &&;
\end{codeblock}

To [expected.object.obs], add:

\begin{itemdescr}
\begin{codeblock}
template<class... Args> constexpr T value_or_construct(Args&&... args) const&;
\end{codeblock}
\pnum
\mandates
\tcode{is_copy_constructible_v<T>} is \tcode{true} and
\tcode{is_constructible_v<T, Args...>} is \tcode{true}.
\pnum
\effects
Equivalent to:
\tcode{return has_value() ? **this : T(std::forward<Args>(args)...);}
\end{itemdescr}

\begin{itemdescr}
\begin{codeblock}
template<class... Args> constexpr T value_or_construct(Args&&... args) &&;
\end{codeblock}
\pnum
\mandates
\tcode{is_move_constructible_v<T>} is \tcode{true} and
\tcode{is_constructible_v<T, Args...>} is \tcode{true}.
\pnum
\effects
Equivalent to:
\tcode{return has_value() ? std::move(**this) : T(std::forward<Args>(args)...);}
\end{itemdescr}

\begin{itemdescr}
\begin{codeblock}
template<class U, class... Args>
  constexpr T value_or_construct(initializer_list<U> il, Args&&... args) const&;
\end{codeblock}
\pnum
\mandates
\tcode{is_copy_constructible_v<T>} is \tcode{true} and
\tcode{is_constructible_v<T, initializer_list<U>\&, Args...>} is \tcode{true}.
\pnum
\effects
Equivalent to:
\tcode{return has_value() ? **this : T(il, std::forward<Args>(args)...);}
\end{itemdescr}

\begin{itemdescr}
\begin{codeblock}
template<class U, class... Args>
  constexpr T value_or_construct(initializer_list<U> il, Args&&... args) &&;
\end{codeblock}
\pnum
\mandates
\tcode{is_move_constructible_v<T>} is \tcode{true} and
\tcode{is_constructible_v<T, initializer_list<U>\&, Args...>} is \tcode{true}.
\pnum
\effects
Equivalent to:
\tcode{return has_value() ? std::move(**this) : T(il, std::forward<Args>(args)...);}
\end{itemdescr}

\begin{itemdescr}
\begin{codeblock}
template<class F> constexpr T value_or_else(F&& f) const&;
\end{codeblock}
\pnum
Let \tcode{U} be \tcode{invoke_result_t<F>}.
\pnum
\mandates
\tcode{is_copy_constructible_v<T>} is \tcode{true} and
\tcode{is_convertible_v<U, T>} is \tcode{true}.
\pnum
\effects
Equivalent to:
\tcode{return has_value() ? **this : T(std::forward<F>(f)());}
\end{itemdescr}

\begin{itemdescr}
\begin{codeblock}
template<class F> constexpr T value_or_else(F&& f) &&;
\end{codeblock}
\pnum
Let \tcode{U} be \tcode{invoke_result_t<F>}.
\pnum
\mandates
\tcode{is_move_constructible_v<T>} is \tcode{true} and
\tcode{is_convertible_v<U, T>} is \tcode{true}.
\pnum
\effects
Equivalent to:
\tcode{return has_value() ? std::move(**this) : T(std::forward<F>(f)());}
\end{itemdescr}

\pnum
\tcode{expected<void,E>} gains neither member --- there is no value to
construct --- and no \tcode{error_or_construct} or \tcode{error_or_else} is
proposed, following \cite{P3413R0}.

\subsection{Reference specializations, and the expected-over-references dependency}

\tcode{std::optional<T\&>} \cite{P2988R12} is in C++26;
\tcode{std::expected<T\&,E>} is not standardized yet, but is expected to
follow. On both, these members return a value rather than a reference ---
exactly as \tcode{value_or} does on those specializations --- and are declared
\tcode{const}, there being no rvalue overload to distinguish. For
\tcode{optional<T\&>} and \tcode{expected<T\&,E>}, add:

\begin{codeblock}
template<class... Args>
  constexpr remove_cv_t<T> value_or_construct(Args&&... args) const;
template<class U, class... Args>
  constexpr remove_cv_t<T> value_or_construct(initializer_list<U> il, Args&&... args) const;
template<class F>
  constexpr remove_cv_t<T> value_or_else(F&& f) const;
\end{codeblock}

\begin{itemdescr}
\pnum
Let \tcode{X} be \tcode{remove_cv_t<T>}, and for \tcode{value_or_else} let
\tcode{U} be \tcode{invoke_result_t<F>}.
\pnum
\mandates
\tcode{is_convertible_v<T\&, X>} is \tcode{true}; for
\tcode{value_or_construct}, \tcode{is_constructible_v<X, Args...>} (respectively
\tcode{is_constructible_v<X, initializer_list<U>\&, Args...>}) is \tcode{true};
for \tcode{value_or_else}, \tcode{is_convertible_v<U, X>} is \tcode{true}.
\pnum
\effects
Equivalent to, respectively:
\begin{codeblock}
return has_value() ? **this : X(std::forward<Args>(args)...);
return has_value() ? **this : X(il, std::forward<Args>(args)...);
return has_value() ? **this : std::forward<F>(f)();
\end{codeblock}
\end{itemdescr}

\pnum
The \tcode{optional<T\&>} wording above is a plain addition to a C++26
facility; nothing about it is contingent. Only the \tcode{expected<T\&,E>}
case is: it depends on the expected-over-references work, which is not yet in
the working draft. Both reference forms are implemented already --- the Beman
\tcode{expected} carries \tcode{expected<T\&,E>}, \tcode{expected<T,E\&>}, and
\tcode{expected<T\&,E\&>}, and the Beman \tcode{optional} the adopted
\tcode{optional<T\&>} --- so both forms run today (see Implementation
Experience).

\section{Feature test macro}

Add a new feature test macro \tcode{__cpp_lib_value_or}, set to the date of
adoption, for the free functions. For the member functions, bump
\tcode{__cpp_lib_optional} and \tcode{__cpp_lib_expected} to the date of
adoption, per \cite{P3413R0}; these are independent of each other and of the
\tcode{optional<T\&>} macro.

\chapter{Impact on the standard}

A pure library extension. The free functions change no existing class or
function. The member functions add \tcode{value_or_construct} and
\tcode{value_or_else} to \tcode{optional} and \tcode{expected} without
altering any existing member, so nothing existing changes meaning; being
member function templates, they add nothing to the ABI. The free functions
require C++23 library support (\tcode{reference_constructs_from_temporary_v});
the members impose no requirement beyond the class they are added to.

\chapter{Acknowledgments}

Thanks to Peter Sommerlad, co-author of \cite{P2988R12}, where the member
\tcode{value_or} question was fought to a standstill, and to the LEWG
reviewers of P1255 and P2988 whose objections shaped these functions. The
Tokyo review homework on \tcode{value_or} alternatives surveyed the
implementation divergence reported here. Thanks to Marc Mutz, whose
\cite{P2218R0} originated the member \tcode{value_or_construct} and
\tcode{value_or_else} that \cite{P3413R0} revived, and to Barry Revzin for the
wording feedback carried through that paper.

\chapter{Document history}

\begin{itemize}
\item \textbf{R1} \today\ --- Unified with \cite{P3413R0} following LEWG's
  Brno review, which asked for a single proposal. Corentin Jabot joins as
  co-author, and the members \tcode{value_or_construct} and
  \tcode{value_or_else} are folded into the proposal, wording, and
  implementation. Independently of the unification: an rvalue nullable moves
  its payload, \tcode{nullable} is proposed as a public concept and stated so
  that const-observability holds, \tcode{reference_or}'s dangling guards are
  corrected and extended to temporary owning wrappers, reference results are
  required where appropriate, \tcode{or_invoke} forwards its invocable, and
  the asks to LEWG are cut from nine to zero. See Changes Since Last Version.
\item \textbf{R0} Brno, 2026-06 --- Initial draft, reviewed by LEWG. The
  free-function follow-on promised in \cite{P2988R12}, concluding the line of
  work begun in \cite{P1255R13}.
\end{itemize}

\renewcommand{\bibname}{References}
\bibliographystyle{abstract}
\bibliography{wg21,mybiblio}

\end{document}
